\documentclass[10pt,a4paper]{amsart}
\usepackage[margin=19mm]{geometry}
\usepackage{amsmath,amssymb,mathtools}
\usepackage[scr=rsfs,cal=euler]{mathalfa}
\usepackage{xfrac}
\usepackage[unicode,hidelinks,bookmarks=false]{hyperref}
\newcommand{\ie}{i.e.\ }
\newcommand{\cf}{cf.\ }
\newcommand{\resp}{respectively\ }
\newcommand{\Subsection}{\subsection}
\providecommand{\nobreakdash}{\nobreak\hbox{-}}
\setlength{\emergencystretch}{2em}
\multlinegap0pt
\usepackage{xcolor}
\usepackage{amssymb}
\usepackage{tikz}
\usepackage[shortlabels]{enumitem}
\newtheorem{lem}{Lemma}
\title{Invariant Sphere Theorem and Ring-Coupled Systems}
\author{Pedro Soares}
\date{Source: arXiv:2608.28223v1 (CC BY 4.0)}
\begin{document}
\maketitle

\noindent\emph{Source and licence.} Invariant Sphere Theorem and Ring-Coupled Systems. Version arXiv:2608.28223v1 (CC BY 4.0), distributed by arXiv under the Creative Commons Attribution 4.0 International licence, http://creativecommons.org/licenses/by/4.0/, as recorded in the arXiv licence metadata for this exact version and shown on the abstract page https://arxiv.org/abs/2608.28223v1. Full licence notice: \texttt{licenses/arxiv-2608.28223-CC-BY-4.0.txt}.\par\smallskip

\noindent\textit{Editorial scope.} Counted unit: the article's lem lem:2dsystem (source0.tex lines 779--791). The article's complete original proof of it is delivered with it (source0.tex lines 793--867, one \texttt{proof} environment of 75 lines). No mathematics is written, repaired or re-derived: the statement and the proof are the article's own lines, in the article's own notation, and no retained line is substituted or re-flowed.  Retained uncounted as the source-local context the proof needs: lines 298--300.  Nothing was omitted from any retained span, so the package carries no omission record.  No retained line contains a citation, so the wrapper carries no bibliography and no bibliography entry.  No retained line contains a figure environment, an image-inclusion command or drawing code, so no image is delivered and the package declares no asset dependency.  Editorial formatting only, in addition: an AMS \texttt{amsart} preamble that restates the article's own declarations and macros used by the retained lines, namely the environments \texttt{lem} on separate counters, together with the standard packages those lines need.  The retained environments print in this document as Lemma 1 in order of appearance, whereas the article prints the same items with the numbering of its own full document.  The complete native source of the article as distributed in the arXiv e-print for this exact version is bundled unchanged as \texttt{sources/208771/source0.tex}.
\begin{equation}\label{eq:edosphere} 
\dot{u}=\mathcal{P}_N(u),
\end{equation}

\begin{lem}\label{lem:2dsystem}
 Let $q(y,z)=a y^3+b y^2 z+c y z^2+d z^3$.	
	The system (\ref{eq:edosphere}) in the half-sphere $\{u\in S^{n-1}:u_1>0 \}$ around the point $e_1\in S^{n-1}$ is equivalent to the following $n-1$ dimensional system:
$$\dot{p}_i=\dfrac{q(p_i,p_{i-1})-p_i q(1,p_n)}{1+p_2^2+\dots +p_n^2},$$
where $p_1=1$, $p_i\in\mathbb{R}$ and $i=2,\dots,n$.

For $n=3$, the Taylor expansion of the previous system around $(p_2,p_3)=(0,0)$ is given by 
$$\begin{cases}
\dot{p}_2 =d+(c-a)p_2+(b-d)p_2^2  -b p_3 p_2-d p_3^2+(2a-c) p_2^3 +(a-2c)p_3^2 p_2 \\
\dot{p}_3 = -a p_3- bp_3^2+a p_3 p_2^2+(2a-c)p_3^3+bp_3^2p_2+cp_3p_2^2+dp_2^3 
\end{cases}+\mathcal{O}(4),$$
where $\mathcal{O}(4)$ denotes the terms of order four or higher in $p_2$ and $p_3$.
\end{lem}

\begin{proof}

We consider the following stereographic projection from the tangent space $T_{e_1}S^{n-1}$ to the half-sphere $\{u\in S^{n-1}:u_1>0 \}$.
Let $P_{e_1}=\{(0,p_2,\dots,p_n)\in \mathbb{R}^n:p_i\in\mathbb{R}\}=T_{e_1}S^{n-1}\cong \mathbb{R}^{n-1}$ and $\phi_{1}:\mathbb{R}^{n-1}\rightarrow S^{n-1}$ given by
$$ (p_2,\dots,p_n)\mapsto \dfrac{(1,p_2,\dots,p_n)}{\|(1,p_2,\dots,p_n)\|}.$$
Note that $\phi_{1}(\mathbb{R}^{n-1})=\{u\in S^{n-1}:u_1>0 \}\subset S^{n-1}$ and 
$$\phi^{-1}_{e_1}(u)=(\dfrac{u_2}{u_1},\dfrac{u_3}{u_1},\dots, \dfrac{u_n}{u_1}).$$
This change of coordinates is invertible and smooth.

Around the equilibrium $e_1$, the system (\ref{eq:edosphere}) can be written in terms of $p=\phi^{-1}_{e_1}(u)\in \mathbb{R}^{n-1}$ as
$$\dot{p}=  D \phi^{-1}_{e_1}(\phi(p)).\dot{u}=D \phi^{-1}_{e_1}(\phi(p)).(Q_N(\phi(p))-\langle Q_N(\phi(p)),\phi(p)\rangle \phi(p))$$
Note that 
$$D \phi^{-1}_{e_1}(u)=\begin{bmatrix}
-\dfrac{u_2}{u_1^2}& \dfrac{1}{u_1}&0&\dots &0\\
-\dfrac{u_3}{u_1^2}& 0&\dfrac{1}{u_1}&\dots &0\\
\vdots& \vdots& \vdots& \ddots& \vdots\\
-\dfrac{u_n}{u_1^2}& 0&0&\dots &\dfrac{1}{u_1}\\
\end{bmatrix}\quad \textrm{ and }\quad D\phi^{-1}_{e_1}(u) u=0.$$
and 
$$\dot{p}= \dfrac{\|(1,p_2,\dots, p_n)\|}{\|(1,p_2,\dots, p_n)\|^3}\begin{bmatrix}
-p_2& 1 &0&\dots &0\\
-p_3& 0&1&\dots &0\\
\vdots& \vdots& \vdots& \ddots& \vdots\\
-p_n & 0&0&\dots &1
\end{bmatrix}Q_N(1,p_2,\dots,p_3)$$
$$=
\dfrac{1}{\|(1,p_2,\dots, p_n)\|^2}\begin{bmatrix}
-p_2& 1 &0&\dots &0\\
-p_3& 0&1&\dots &0\\
\vdots& \vdots& \vdots& \ddots& \vdots\\
-p_n & 0&0&\dots &1
\end{bmatrix}\begin{bmatrix}
q(1,p_n)\\
q(p_2,1)\\
\vdots\\
q(p_n,p_{n-1})
\end{bmatrix}
$$
$$
=\dfrac{1}{1+p_2^2+\dots +p_n^2}\begin{bmatrix}
q(p_2,1)-p_2q(1,p_n)\\
q(p_3,p_2)-p_3q(1,p_n)\\
\vdots\\
q(p_n,p_{n-1})-p_n q(1,p_n)\\
\end{bmatrix}.$$
So $$\dot{p}_i=\dfrac{q(p_i,p_{i-1})-p_i q(1,p_n)}{1+p_2^2+\dots +p_n^2},$$
where $p_1=1$ and $i=2,\dots,n$.

We have the following Taylor expansion up to order $4$ of 
$$\dfrac{1}{(1+p_2^2+\dots + p_n^2)}=1-(p_2^2+\dots + p_n^2)+\mathcal{O}(4)$$
So the vector field can be written as 
$$\dfrac{q(p_i,p_{i-1})-p_i q(1,p_n)}{1+p_2^2+\dots +p_n^2}=q(p_i,p_{i-1})-p_i q(1,p_n)-(p_2^2+\dots +p_n^2)(q(p_i,p_{i-1})-p_i q(1,p_n))+\mathcal{O}(4)$$

For $i=2$, we have 
\begin{align*}
\dot{p}_2=&q(p_2,p_1)-p_2 q(1,p_n)-(p_2^2+\dots +p_n^2)(q(p_2,p_1)-p_2 q(1,p_n))+\mathcal{O}(4)\\
 =&d+(c-a)p_2+(b-d)p_2^2 -d(p_3^2+\dots +p_n^2) -b p_n p_2\\
 & +(2a-c) p_2^3 +(a-c) p_2(p_3^2+\dots +p_{n-1}^2) +(a-2c)p_n^2 p_2 +\mathcal{O}(4)
\end{align*}
For $i=3,\dots,n$, we have 
\begin{align*}
\dot{p}_i	=&q(p_i,p_{i-1})-p_i q(1,p_n)-(p_2^2+\dots +p_n^2)(q(p_i,p_{i-1})-p_i q(1,p_n))+\mathcal{O}(4)\\
					=& -ap_i -bp_np_i\\
					&+q(p_i,p_{i-1})+a p_i(p_2^2+\dots +p_{n-1}^2)+(a-c)p_i p_n^2 +\mathcal{O}(4) 
\end{align*}

For $n=3$, the system (\ref{eq:edosphere}) in the half-sphere around the point $(1,0,0)\in S^2$ is equivalent to the following two-dimensional system:
$$\begin{cases}
\dot{p}_2=d+(c-a)p_2+(b-d)p_2^2  -b p_3 p_2-d p_3^2+(2a-c) p_2^3 +(a-2c)p_3^2 p_2 \\
\dot{p}_3	= -a p_3- bp_3^2+a p_3 p_2^2+(2a-c)p_3^3+bp_3^2p_2+cp_3p_2^2+dp_2^3 
\end{cases}+\mathcal{O}(4)
$$


\end{proof}

\end{document}
